
% ============================================================
% Research Track: Triadic gravity-response diagnostics
% Derived from the EM--gravity sector-split conversation track.
% ============================================================

\section{Research Track: Flame-Lock, Photonless-Caustic, and Elastic-Shell Diagnostics}
\label{sec:rt_flame_photonless_shell}

\paragraph{Status.}
\textbf{[RESEARCH TRACK / DIAGNOSTIC ONLY].}
This section records a three-channel diagnostic motivated by anomalous-field
witness reports. The reports are not used as evidence. They serve only as a
catalogue of separable experimental signatures compatible with the canonical
sector-projection law of Section~\ref{subsec:canonical_em_gravity_closure}.

\subsection{Mode I: flame-lock / plume stabilization}

A flame or thermal plume may be stabilized without significant optical ray
deflection if the dominant SST response is local impedance or damping rather
than photon-index curvature. A minimal diagnostic equation is
\begin{align}
    \frac{\partial\mathbf{u}_{\mathrm{plume}}}{\partial t}
    +
    (\mathbf{u}_{\mathrm{plume}}\cdot\nabla)\mathbf{u}_{\mathrm{plume}}
    =
    -\frac{1}{\rho_{\mathrm{plume}}}\nabla p
    +
    \nu\nabla^2\mathbf{u}_{\mathrm{plume}}
    -
    \Gamma_{\swirlarrow}\mathbf{u}_{\mathrm{plume}}.
    \label{eq:rt_flame_lock_plume_equation}
\end{align}
The diagnostic signature is
\begin{align}
    P_{\mathrm{plume}}(f)\downarrow,
    \qquad
    \Delta\phi_\gamma\approx0,
    \qquad
    \nabla_\perp n_\gamma\approx0.
    \label{eq:rt_flame_lock_signature}
\end{align}
This channel is a boundary/impedance response, not a standalone proof of a
gravitational bulk force.

\textbf{Primary orthodox confounders:}
ion wind, acoustic forcing, thermal convection suppression, electric-field
stabilization of charged flame species, camera exposure artifacts, and ordinary
airflow.

\subsection{Mode II: photonless sphere / optical caustic}

A dark or photon-depleted region is interpreted, if real, as an optical caustic,
shadow, or phase-gradient effect, not as a literal astrophysical black hole. The
minimal ray equation is
\begin{align}
    \frac{d}{ds}
    \left(
    n_\gamma\frac{d\mathbf{x}}{ds}
    \right)
    =
    \nabla n_\gamma,
    \qquad
    n_\gamma=\SwirlClock^{-1}.
    \label{eq:rt_photonless_caustic_ray}
\end{align}
The diagnostic signature is
\begin{align}
    \nabla_\perp n_\gamma\neq0,
    \qquad
    \Delta\phi_\gamma
    =
    \frac{2\pi}{\lambda}
    \int(n_\gamma-1)\,ds
    \neq0,
    \qquad
    \mathbf{f}_{\mathrm{matter}}\approx0.
    \label{eq:rt_photonless_caustic_signature}
\end{align}
Macroscopic photon depletion or black-ball-like appearance over laboratory path
lengths requires a resonant enhancement factor:
\begin{align}
    n_\gamma-1
    =
    Q_\gamma
    \frac{\lVert\mathbf{u}_{\swirlarrow}\rVert^2}{2c^2},
    \qquad
    Q_\gamma\gg1.
    \label{eq:rt_optical_gain_factor}
\end{align}
The gain \(Q_\gamma\) is apparatus-specific and is not a canonical constant.

\textbf{Primary orthodox confounders:}
thermal lensing, plasma absorption, smoke or aerosol scattering, detector
saturation, schlieren artifacts, camera auto-exposure, and ordinary optical
shadowing.

\subsection{Mode III: elastic shell / boundary-stress response}

A localized repulsive shell is represented as a boundary-stress potential,
\begin{align}
    U_{\mathrm{shell}}(r)
    =
    U_0
    \exp\left[
    -\frac{(r-R_s)^2}{2\sigma_s^2}
    \right],
    \label{eq:rt_shell_potential}
\end{align}
with radial force
\begin{align}
    F_r(r)
    =
    -\frac{dU_{\mathrm{shell}}}{dr}
    =
    \frac{U_0(r-R_s)}{\sigma_s^2}
    \exp\left[
    -\frac{(r-R_s)^2}{2\sigma_s^2}
    \right].
    \label{eq:rt_shell_force}
\end{align}
The diagnostic signature is
\begin{align}
    \nabla\cdot\boldsymbol{\sigma}^{\mathrm{swirl}}\neq0,
    \qquad
    \Delta\phi_\gamma\approx0
    \quad
    \text{or independently measurable}.
    \label{eq:rt_shell_signature}
\end{align}
This mode is a stress/boundary response. It must not be rebranded as gravity
shielding unless the bulk acceleration sector is independently measured.

\textbf{Primary orthodox confounders:}
electrostatic repulsion, magnetic force, acoustic radiation pressure, airflow,
vibration, thermal expansion, hidden mechanical coupling, and contact charging.

\subsection{Required measurement vector}

A valid test must record the full sectoral observable set
\begin{align}
    \boxed{
    \left\{
    F(t),
    \Delta\phi_\gamma(t),
    \frac{\Delta f}{f}(t),
    P_{\mathrm{plume}}(f,t),
    \mathbf{E}(t),
    \mathbf{B}(t),
    T(t),
    p_{\mathrm{air}}(t)
    \right\}.
    }
    \label{eq:rt_triadic_measurement_vector}
\end{align}
Any reported force must first be decomposed as
\begin{align}
    \mathbf{F}_{\mathrm{meas}}
    =
    \mathbf{F}_{\mathrm{EM}}
    +
    \mathbf{F}_{\mathrm{EHD}}
    +
    \mathbf{F}_{\mathrm{thermal}}
    +
    \mathbf{F}_{\mathrm{acoustic}}
    +
    \mathbf{F}_{\mathrm{mechanical}}
    +
    \mathbf{F}_{\omega}
    +
    \mathbf{F}_{\mathrm{res}}.
    \label{eq:rt_triadic_force_decomposition}
\end{align}
Only \(\mathbf{F}_{\mathrm{res}}\), after uncertainty propagation and null tests,
is admissible as an SST research-track candidate.

\subsection{Null hierarchy}

The minimum null hierarchy is:
\begin{enumerate}
    \item reproduce the measurement with all high-voltage and high-current channels disabled;
    \item randomize phase order at fixed RMS power;
    \item replace the structured geometry with a symmetry-preserving dummy load of matched impedance;
    \item repeat in still air, shielded enclosure, and reduced pressure where safe;
    \item compare optical, force, and plume observables under the same drive waveform.
\end{enumerate}
A positive SST candidate requires a correlated sectoral signature that survives
these nulls and scales with the predicted mode-selection or stress-closure
parameter.

\section{Research Track: Atomic Gravity Closure and Condensate Coherence}
\label{sec:rt_atomic_gravity_closure_condensate}

\paragraph{Status.}
\textbf{[CANON-CANDIDATE / RESEARCH BRIDGE].}
This section records the electron-envelope and condensate-coherence clue for
future gravity-sector tests. It does not claim gravitational shielding.

\subsection{Electron envelope as atomic closure layer}

In a neutral atom, the electron envelope cancels the far electromagnetic charge
sector but does not cancel positive swirl-energy density. Therefore the electron
is not the dominant mass reservoir, but it is the outer phase and linking
boundary through which the neutral atom closes its electromagnetic projection:
\begin{align}
    Q_{\mathrm{atom}}=0,
    \qquad
    \rhoM
    =
    \frac{\rhoE}{c^2}
    >0.
    \label{eq:rt_atomic_charge_mass_split}
\end{align}
The canon-safe interpretation is
\begin{equation}
\boxed{
\text{the nucleus supplies most mass-energy, while the electron envelope supplies
      the neutral phase boundary and long-range matching layer.}
}
\label{eq:rt_atomic_electron_envelope_summary}
\end{equation}

\subsection{Condensate coherence extension}

For coherent electronic or atomic condensates, a macroscopic order parameter
\begin{align}
    \Psi_{\mathrm{cond}}(\mathbf{x})
    =
    \sqrt{n_s(\mathbf{x})}\,e^{i\theta(\mathbf{x})}
    \label{eq:rt_condensate_order_parameter}
\end{align}
may coherently reweight boundary-stress and optical/clock projections:
\begin{align}
    \boldsymbol{\sigma}^{\mathrm{swirl}}
    &\rightarrow
    \boldsymbol{\sigma}^{\mathrm{swirl}}_{\mathrm{coh}}[n_s,\theta,\nabla\theta],
    \\
    n_\gamma
    &\rightarrow
    n_{\gamma,\mathrm{coh}}.
    \label{eq:rt_condensate_projection_reweighting}
\end{align}
The proposed clue from Meissner/BEC physics is therefore not mass cancellation,
but collective phase closure: many microscopic carriers can impose one
macroscopic boundary condition.

\textbf{[CRITICAL NOTE].}
This does not imply gravitational shielding or mass cancellation. The testable
claim is only a possible phase-dependent change in boundary-stress, optical
phase, or clock response across a coherence transition such as \(T_c\).

\subsection{Condensate test signature}

A minimally admissible superconducting or BEC test compares the same apparatus
above and below the coherence transition:
\begin{align}
    \Delta\mathcal{O}_{\mathrm{coh}}
    =
    \mathcal{O}(T<T_c)-\mathcal{O}(T>T_c),
    \qquad
    \mathcal{O}
    \in
    \left\{
    F,
    \Delta\phi_\gamma,
    \Delta f/f,
    \Delta\Phi
    \right\}.
    \label{eq:rt_condensate_test_signature}
\end{align}
A positive result must track coherence and vanish under matched non-coherent
thermal, electromagnetic, and mechanical controls.
